\documentclass[aspectratio=169,11pt]{beamer}
\usetheme{Madrid}
\usecolortheme{default}
\setbeamertemplate{navigation symbols}{}
\setbeamertemplate{footline}[frame number]
\usepackage{enumitem}
\usepackage{amsmath,amssymb}
\usepackage{booktabs}
\usepackage{array}

\title{Databricks ML Associate}
\subtitle{Questions 201--300 (Third 100)}
\author{Practice Exam Deck}
\date{}

\begin{document}
	
	\begin{frame}
		\titlepage
	\end{frame}
	
	\begin{frame}{How to Use This Deck}
		\begin{itemize}
			\item 100 new questions (Q201--Q300), none repeated from the first two sets.
			\item Each question slide has four options (A--D).
			\item The next slide gives the correct answer and explains why the right option is right and why each wrong option is wrong.
		\end{itemize}
	\end{frame}
	
	% ---------------------------------------------------------------
	\begin{frame}{Question 201}
		Which method updates the description of an existing Feature Store table?
		\begin{enumerate}[label=\Alph*.]
			\item \texttt{fe.update\_table()}
			\item \texttt{fe.set\_description()}
			\item \texttt{fe.write\_table()}
			\item \texttt{fe.create\_table()}
		\end{enumerate}
	\end{frame}
	\begin{frame}{Answer 201}
		\textbf{Correct: A} --- \texttt{update\_table()} updates table metadata.\\
		\vspace{0.5em}
		\textbf{Wrong:}
		\begin{itemize}
			\item B is not the API.
			\item C writes features.
			\item D creates a table.
		\end{itemize}
	\end{frame}
	
	% ---------------------------------------------------------------
	\begin{frame}{Question 202}
		Which method deletes a feature table?
		\begin{enumerate}[label=\Alph*.]
			\item \texttt{fe.drop\_table()}
			\item \texttt{fe.delete\_table()}
			\item \texttt{fe.remove\_table()}
			\item \texttt{fe.write\_table()}
		\end{enumerate}
	\end{frame}
	\begin{frame}{Answer 202}
		\textbf{Correct: A} --- \texttt{drop\_table()} deletes a feature table.\\
		\vspace{0.5em}
		\textbf{Wrong:}
		\begin{itemize}
			\item B, C are not the API.
			\item D writes features.
		\end{itemize}
	\end{frame}
	
	% ---------------------------------------------------------------
	\begin{frame}{Question 203}
		What happens if two feature table rows share the same primary key without a timestamp key?
		\begin{enumerate}[label=\Alph*.]
			\item The table cannot be used for point-in-time lookups
			\item It trains faster
			\item It improves accuracy
			\item It creates a new feature table
		\end{enumerate}
	\end{frame}
	\begin{frame}{Answer 203}
		\textbf{Correct: A} --- point-in-time lookups require a timestamp key.\\
		\vspace{0.5em}
		\textbf{Wrong:}
		\begin{itemize}
			\item B, C are false.
			\item D does not happen.
		\end{itemize}
	\end{frame}
	
	% ---------------------------------------------------------------
	\begin{frame}{Question 204}
		Which is a valid reason to use a Feature Store over raw Delta tables?
		\begin{enumerate}[label=\Alph*.]
			\item Automated point-in-time joins and feature lineage
			\item Faster cluster startup
			\item Automatic model deployment
			\item Free storage
		\end{enumerate}
	\end{frame}
	\begin{frame}{Answer 204}
		\textbf{Correct: A} --- Feature Store provides point-in-time joins and lineage.\\
		\vspace{0.5em}
		\textbf{Wrong:}
		\begin{itemize}
			\item B is compute.
			\item C is deployment.
			\item D is billing.
		\end{itemize}
	\end{frame}
	
	% ---------------------------------------------------------------
	\begin{frame}{Question 205}
		What does \texttt{fs.get\_table()} return?
		\begin{enumerate}[label=\Alph*.]
			\item A FeatureStoreTable object
			\item A Spark DataFrame
			\item A pandas DataFrame
			\item A model
		\end{enumerate}
	\end{frame}
	\begin{frame}{Answer 205}
		\textbf{Correct: A} --- it returns a FeatureStoreTable object.\\
		\vspace{0.5em}
		\textbf{Wrong:}
		\begin{itemize}
			\item B you call \texttt{.load\_df()} for a DataFrame.
			\item C is pandas.
			\item D is a model.
		\end{itemize}
	\end{frame}
	
	% ---------------------------------------------------------------
	\begin{frame}{Question 206}
		Which MLflow component stores model files?
		\begin{enumerate}[label=\Alph*.]
			\item The artifact store
			\item The backend store
			\item The tracking server
			\item The model registry
		\end{enumerate}
	\end{frame}
	\begin{frame}{Answer 206}
		\textbf{Correct: A} --- model files are artifacts.\\
		\vspace{0.5em}
		\textbf{Wrong:}
		\begin{itemize}
			\item B stores metadata.
			\item C coordinates tracking.
			\item D manages versions.
		\end{itemize}
	\end{frame}
	
	% ---------------------------------------------------------------
	\begin{frame}{Question 207}
		What does \texttt{mlflow.log\_dict()} do?
		\begin{enumerate}[label=\Alph*.]
			\item Logs a dictionary as a JSON artifact
			\item Logs a parameter
			\item Logs a metric
			\item Logs a tag
		\end{enumerate}
	\end{frame}
	\begin{frame}{Answer 207}
		\textbf{Correct: A} --- \texttt{log\_dict()} stores a JSON artifact.\\
		\vspace{0.5em}
		\textbf{Wrong:}
		\begin{itemize}
			\item B is \texttt{log\_param}.
			\item C is \texttt{log\_metric}.
			\item D is \texttt{set\_tag}.
		\end{itemize}
	\end{frame}
	
	% ---------------------------------------------------------------
	\begin{frame}{Question 208}
		What is an MLflow run?
		\begin{enumerate}[label=\Alph*.]
			\item A single execution of training code
			\item A cluster
			\item A feature table
			\item An experiment
		\end{enumerate}
	\end{frame}
	\begin{frame}{Answer 208}
		\textbf{Correct: A} --- a run is one execution.\\
		\vspace{0.5em}
		\textbf{Wrong:}
		\begin{itemize}
			\item B is compute.
			\item C is data.
			\item D groups runs.
		\end{itemize}
	\end{frame}
	
	% ---------------------------------------------------------------
	\begin{frame}{Question 209}
		What is an MLflow experiment?
		\begin{enumerate}[label=\Alph*.]
			\item A collection of related runs
			\item A single run
			\item A model
			\item A cluster
		\end{enumerate}
	\end{frame}
	\begin{frame}{Answer 209}
		\textbf{Correct: A} --- experiments group related runs.\\
		\vspace{0.5em}
		\textbf{Wrong:}
		\begin{itemize}
			\item B is a run.
			\item C is a model.
			\item D is compute.
		\end{itemize}
	\end{frame}
	
	% ---------------------------------------------------------------
	\begin{frame}{Question 210}
		How can you compare runs visually in MLflow?
		\begin{enumerate}[label=\Alph*.]
			\item Use the experiment UI to compare metrics and parameters
			\item Open a notebook
			\item Use a SQL warehouse
			\item Use a cluster log
		\end{enumerate}
	\end{frame}
	\begin{frame}{Answer 210}
		\textbf{Correct: A} --- the experiment UI compares runs.\\
		\vspace{0.5em}
		\textbf{Wrong:}
		\begin{itemize}
			\item B, C, D do not compare runs.
		\end{itemize}
	\end{frame}
	
	% ---------------------------------------------------------------
	\begin{frame}{Question 211}
		Which MLflow API loads a model from a registry URI?
		\begin{enumerate}[label=\Alph*.]
			\item \texttt{mlflow.pyfunc.load\_model()}
			\item \texttt{mlflow.load\_model()}
			\item \texttt{mlflow.get\_model()}
			\item \texttt{mlflow.fetch\_model()}
		\end{enumerate}
	\end{frame}
	\begin{frame}{Answer 211}
		\textbf{Correct: A} --- \texttt{pyfunc.load\_model()} loads a registry model.\\
		\vspace{0.5em}
		\textbf{Wrong:}
		\begin{itemize}
			\item B, C, D are not the loading API.
		\end{itemize}
	\end{frame}
	
	% ---------------------------------------------------------------
	\begin{frame}{Question 212}
		What is MLflow autologging?
		\begin{enumerate}[label=\Alph*.]
			\item Automatic logging of params, metrics, and models
			\item Automatic model deployment
			\item Automatic data cleaning
			\item Automatic cluster creation
		\end{enumerate}
	\end{frame}
	\begin{frame}{Answer 212}
		\textbf{Correct: A} --- autologging logs run data automatically.\\
		\vspace{0.5em}
		\textbf{Wrong:}
		\begin{itemize}
			\item B is deployment.
			\item C is data prep.
			\item D is compute.
		\end{itemize}
	\end{frame}
	
	% ---------------------------------------------------------------
	\begin{frame}{Question 213}
		Which MLflow function enables autologging for scikit-learn?
		\begin{enumerate}[label=\Alph*.]
			\item \texttt{mlflow.sklearn.autolog()}
			\item \texttt{mlflow.autolog\_sklearn()}
			\item \texttt{mlflow.enable\_sklearn()}
			\item \texttt{mlflow.sklearn.log()}
		\end{enumerate}
	\end{frame}
	\begin{frame}{Answer 213}
		\textbf{Correct: A} --- \texttt{sklearn.autolog()} enables autologging.\\
		\vspace{0.5em}
		\textbf{Wrong:}
		\begin{itemize}
			\item B, C, D are not the API.
		\end{itemize}
	\end{frame}
	
	% ---------------------------------------------------------------
	\begin{frame}{Question 214}
		Which MLflow API creates a new experiment?
		\begin{enumerate}[label=\Alph*.]
			\item \texttt{mlflow.create\_experiment()}
			\item \texttt{mlflow.new\_experiment()}
			\item \texttt{mlflow.start\_experiment()}
			\item \texttt{mlflow.set\_experiment()}
		\end{enumerate}
	\end{frame}
	\begin{frame}{Answer 214}
		\textbf{Correct: A} --- \texttt{create\_experiment()} creates one.\\
		\vspace{0.5em}
		\textbf{Wrong:}
		\begin{itemize}
			\item B, C are not the API.
			\item D sets the active experiment.
		\end{itemize}
	\end{frame}
	
	% ---------------------------------------------------------------
	\begin{frame}{Question 215}
		Which MLflow API sets the active experiment?
		\begin{enumerate}[label=\Alph*.]
			\item \texttt{mlflow.set\_experiment()}
			\item \texttt{mlflow.create\_experiment()}
			\item \texttt{mlflow.start\_experiment()}
			\item \texttt{mlflow.use\_experiment()}
		\end{enumerate}
	\end{frame}
	\begin{frame}{Answer 215}
		\textbf{Correct: A} --- \texttt{set\_experiment()} sets the active one.\\
		\vspace{0.5em}
		\textbf{Wrong:}
		\begin{itemize}
			\item B creates.
			\item C, D are not the API.
		\end{itemize}
	\end{frame}
	
	% ---------------------------------------------------------------
	\begin{frame}{Question 216}
		What does AutoML generate in addition to the leaderboard?
		\begin{enumerate}[label=\Alph*.]
			\item A data exploration notebook
			\item A cluster
			\item A SQL warehouse
			\item A Git repo
		\end{enumerate}
	\end{frame}
	\begin{frame}{Answer 216}
		\textbf{Correct: A} --- AutoML generates an EDA notebook.\\
		\vspace{0.5em}
		\textbf{Wrong:}
		\begin{itemize}
			\item B, C, D are not produced by AutoML.
		\end{itemize}
	\end{frame}
	
	% ---------------------------------------------------------------
	\begin{frame}{Question 217}
		What input does AutoML require?
		\begin{enumerate}[label=\Alph*.]
			\item A table or DataFrame with a target column
			\item A trained model
			\item A cluster
			\item A SQL warehouse
		\end{enumerate}
	\end{frame}
	\begin{frame}{Answer 217}
		\textbf{Correct: A} --- AutoML needs data with a target.\\
		\vspace{0.5em}
		\textbf{Wrong:}
		\begin{itemize}
			\item B is the output.
			\item C is compute.
			\item D is SQL compute.
		\end{itemize}
	\end{frame}
	
	% ---------------------------------------------------------------
	\begin{frame}{Question 218}
		Which is a supported AutoML problem type?
		\begin{enumerate}[label=\Alph*.]
			\item Classification
			\item Clustering
			\item Dimensionality reduction
			\item Recommendation
		\end{enumerate}
	\end{frame}
	\begin{frame}{Answer 218}
		\textbf{Correct: A} --- AutoML supports classification (and regression/forecasting).\\
		\vspace{0.5em}
		\textbf{Wrong:}
		\begin{itemize}
			\item B, C, D are not AutoML problem types.
		\end{itemize}
	\end{frame}
	
	% ---------------------------------------------------------------
	\begin{frame}{Question 219}
		What does AutoML do with generated notebooks?
		\begin{enumerate}[label=\Alph*.]
			\item It lets you edit and reproduce the trials
			\item It deletes them
			\item It hides them
			\item It converts them to SQL
		\end{enumerate}
	\end{frame}
	\begin{frame}{Answer 219}
		\textbf{Correct: A} --- generated notebooks are editable.\\
		\vspace{0.5em}
		\textbf{Wrong:}
		\begin{itemize}
			\item B, C, D are false.
		\end{itemize}
	\end{frame}
	
	% ---------------------------------------------------------------
	\begin{frame}{Question 220}
		Which API registers an AutoML best model into the registry?
		\begin{enumerate}[label=\Alph*.]
			\item One-click registration in the UI or \texttt{mlflow.register\_model()}
			\item \texttt{automl.deploy()}
			\item \texttt{automl.register()}
			\item \texttt{mlflow.save()}
		\end{enumerate}
	\end{frame}
	\begin{frame}{Answer 220}
		\textbf{Correct: A} --- registration is one-click or via \texttt{register\_model}.\\
		\vspace{0.5em}
		\textbf{Wrong:}
		\begin{itemize}
			\item B, C are not APIs.
			\item D is not registration.
		\end{itemize}
	\end{frame}
	
	% ---------------------------------------------------------------
	\begin{frame}{Question 221}
		What does \texttt{df.withColumnRenamed()} do?
		\begin{enumerate}[label=\Alph*.]
			\item Renames an existing column
			\item Adds a column
			\item Drops a column
			\item Selects columns
		\end{enumerate}
	\end{frame}
	\begin{frame}{Answer 221}
		\textbf{Correct: A} --- it renames a column.\\
		\vspace{0.5em}
		\textbf{Wrong:}
		\begin{itemize}
			\item B is \texttt{withColumn}.
			\item C is \texttt{drop}.
			\item D is \texttt{select}.
		\end{itemize}
	\end{frame}
	
	% ---------------------------------------------------------------
	\begin{frame}{Question 222}
		What does \texttt{df.drop()} do?
		\begin{enumerate}[label=\Alph*.]
			\item Drops columns
			\item Drops rows
			\item Drops nulls
			\item Drops the DataFrame
		\end{enumerate}
	\end{frame}
	\begin{frame}{Answer 222}
		\textbf{Correct: A} --- \texttt{drop()} removes columns.\\
		\vspace{0.5em}
		\textbf{Wrong:}
		\begin{itemize}
			\item B is \texttt{filter}.
			\item C is \texttt{na.drop}.
			\item D is not a Spark operation.
		\end{itemize}
	\end{frame}
	
	% ---------------------------------------------------------------
	\begin{frame}{Question 223}
		Which Spark function returns the number of rows?
		\begin{enumerate}[label=\Alph*.]
			\item \texttt{count()}
			\item \texttt{size()}
			\item \texttt{length()}
			\item \texttt{numRows()}
		\end{enumerate}
	\end{frame}
	\begin{frame}{Answer 223}
		\textbf{Correct: A} --- \texttt{count()} returns row count.\\
		\vspace{0.5em}
		\textbf{Wrong:}
		\begin{itemize}
			\item B, C are for arrays/strings.
			\item D is not an API.
		\end{itemize}
	\end{frame}
	
	% ---------------------------------------------------------------
	\begin{frame}{Question 224}
		Which operation sorts a Spark DataFrame?
		\begin{enumerate}[label=\Alph*.]
			\item \texttt{orderBy()}
			\item \texttt{sortBy()}
			\item \texttt{arrange()}
			\item \texttt{rank()}
		\end{enumerate}
	\end{frame}
	\begin{frame}{Answer 224}
		\textbf{Correct: A} --- \texttt{orderBy()} sorts rows.\\
		\vspace{0.5em}
		\textbf{Wrong:}
		\begin{itemize}
			\item B is Scala collections.
			\item C is pandas.
			\item D is a window function.
		\end{itemize}
	\end{frame}
	
	% ---------------------------------------------------------------
	\begin{frame}{Question 225}
		What does \texttt{df.limit(10)} return?
		\begin{enumerate}[label=\Alph*.]
			\item A new DataFrame with at most 10 rows
			\item The 10th row
			\item A scalar
			\item A column
		\end{enumerate}
	\end{frame}
	\begin{frame}{Answer 225}
		\textbf{Correct: A} --- it returns up to 10 rows.\\
		\vspace{0.5em}
		\textbf{Wrong:}
		\begin{itemize}
			\item B, C, D are not what it returns.
		\end{itemize}
	\end{frame}
	
	% ---------------------------------------------------------------
	\begin{frame}{Question 226}
		Why is \texttt{collect()} dangerous on large DataFrames?
		\begin{enumerate}[label=\Alph*.]
			\item It pulls all data to the driver, risking memory overflow
			\item It deletes data
			\item It changes the schema
			\item It drops nulls
		\end{enumerate}
	\end{frame}
	\begin{frame}{Answer 226}
		\textbf{Correct: A} --- \texttt{collect()} moves data to the driver.\\
		\vspace{0.5em}
		\textbf{Wrong:}
		\begin{itemize}
			\item B, C, D are false.
		\end{itemize}
	\end{frame}
	
	% ---------------------------------------------------------------
	\begin{frame}{Question 227}
		Which action triggers computation on a Spark DataFrame?
		\begin{enumerate}[label=\Alph*.]
			\item \texttt{show()}
			\item \texttt{select()}
			\item \texttt{filter()}
			\item \texttt{withColumn()}
		\end{enumerate}
	\end{frame}
	\begin{frame}{Answer 227}
		\textbf{Correct: A} --- \texttt{show()} is an action.\\
		\vspace{0.5em}
		\textbf{Wrong:}
		\begin{itemize}
			\item B, C, D are lazy transformations.
		\end{itemize}
	\end{frame}
	
	% ---------------------------------------------------------------
	\begin{frame}{Question 228}
		What is lazy evaluation in Spark?
		\begin{enumerate}[label=\Alph*.]
			\item Transformations are not executed until an action is called
			\item Data is never loaded
			\item Results are cached automatically
			\item Queries run immediately
		\end{enumerate}
	\end{frame}
	\begin{frame}{Answer 228}
		\textbf{Correct: A} --- Spark delays execution until an action.\\
		\vspace{0.5em}
		\textbf{Wrong:}
		\begin{itemize}
			\item B, C, D are false.
		\end{itemize}
	\end{frame}
	
	% ---------------------------------------------------------------
	\begin{frame}{Question 229}
		What is a broadcast join?
		\begin{enumerate}[label=\Alph*.]
			\item A join where a small table is sent to all workers
			\item A join on all columns
			\item A join that drops rows
			\item A join without a key
		\end{enumerate}
	\end{frame}
	\begin{frame}{Answer 229}
		\textbf{Correct: A} --- broadcast joins replicate a small table.\\
		\vspace{0.5em}
		\textbf{Wrong:}
		\begin{itemize}
			\item B, C, D are false.
		\end{itemize}
	\end{frame}
	
	% ---------------------------------------------------------------
	\begin{frame}{Question 230}
		What is a shuffle in Spark?
		\begin{enumerate}[label=\Alph*.]
			\item Redistribution of data across partitions
			\item A random sample
			\item A cache operation
			\item A schema change
		\end{enumerate}
	\end{frame}
	\begin{frame}{Answer 230}
		\textbf{Correct: A} --- shuffles redistribute data.\\
		\vspace{0.5em}
		\textbf{Wrong:}
		\begin{itemize}
			\item B, C, D are false.
		\end{itemize}
	\end{frame}
	
	% ---------------------------------------------------------------
	\begin{frame}{Question 231}
		What is a UDF in Spark?
		\begin{enumerate}[label=\Alph*.]
			\item A user-defined function applied to columns
			\item A built-in function
			\item A DataFrame
			\item A cluster
		\end{enumerate}
	\end{frame}
	\begin{frame}{Answer 231}
		\textbf{Correct: A} --- a UDF is a user-defined function.\\
		\vspace{0.5em}
		\textbf{Wrong:}
		\begin{itemize}
			\item B is built-in.
			\item C is a DataFrame.
			\item D is compute.
		\end{itemize}
	\end{frame}
	
	% ---------------------------------------------------------------
	\begin{frame}{Question 232}
		Why can Python UDFs be slower than native Spark functions?
		\begin{enumerate}[label=\Alph*.]
			\item They require serialization between JVM and Python
			\item They are always parallel
			\item They use no memory
			\item They are written in Scala
		\end{enumerate}
	\end{frame}
	\begin{frame}{Answer 232}
		\textbf{Correct: A} --- Python UDFs pay serialization overhead.\\
		\vspace{0.5em}
		\textbf{Wrong:}
		\begin{itemize}
			\item B, C, D are false.
		\end{itemize}
	\end{frame}
	
	% ---------------------------------------------------------------
	\begin{frame}{Question 233}
		What is a pandas UDF?
		\begin{enumerate}[label=\Alph*.]
			\item A vectorized UDF that operates on pandas Series/DataFrames
			\item A built-in function
			\item A clustering method
			\item A storage format
		\end{enumerate}
	\end{frame}
	\begin{frame}{Answer 233}
		\textbf{Correct: A} --- pandas UDFs are vectorized.\\
		\vspace{0.5em}
		\textbf{Wrong:}
		\begin{itemize}
			\item B, C, D are false.
		\end{itemize}
	\end{frame}
	
	% ---------------------------------------------------------------
	\begin{frame}{Question 234}
		What does \texttt{applyInPandas} require?
		\begin{enumerate}[label=\Alph*.]
			\item A schema for the output DataFrame
			\item A cluster
			\item A model
			\item A feature table
		\end{enumerate}
	\end{frame}
	\begin{frame}{Answer 234}
		\textbf{Correct: A} --- it requires an output schema.\\
		\vspace{0.5em}
		\textbf{Wrong:}
		\begin{itemize}
			\item B, C, D are not required.
		\end{itemize}
	\end{frame}
	
	% ---------------------------------------------------------------
	\begin{frame}{Question 235}
		Which Spark ML evaluator is used for regression?
		\begin{enumerate}[label=\Alph*.]
			\item \texttt{RegressionEvaluator}
			\item \texttt{BinaryClassificationEvaluator}
			\item \texttt{MulticlassClassificationEvaluator}
			\item \texttt{ClusteringEvaluator}
		\end{enumerate}
	\end{frame}
	\begin{frame}{Answer 235}
		\textbf{Correct: A} --- it evaluates regression.\\
		\vspace{0.5em}
		\textbf{Wrong:}
		\begin{itemize}
			\item B, C are classification.
			\item D is clustering.
		\end{itemize}
	\end{frame}
	
	% ---------------------------------------------------------------
	\begin{frame}{Question 236}
		Which metric can \texttt{RegressionEvaluator} compute?
		\begin{enumerate}[label=\Alph*.]
			\item RMSE
			\item AUC
			\item F1
			\item Silhouette
		\end{enumerate}
	\end{frame}
	\begin{frame}{Answer 236}
		\textbf{Correct: A} --- it computes RMSE and related regression metrics.\\
		\vspace{0.5em}
		\textbf{Wrong:}
		\begin{itemize}
			\item B is AUC.
			\item C is F1.
			\item D is clustering.
		\end{itemize}
	\end{frame}
	
	% ---------------------------------------------------------------
	\begin{frame}{Question 237}
		What does \texttt{setMetricName()} do on an evaluator?
		\begin{enumerate}[label=\Alph*.]
			\item Sets the metric to compute
			\item Sets the label column
			\item Sets the prediction column
			\item Sets the model
		\end{enumerate}
	\end{frame}
	\begin{frame}{Answer 237}
		\textbf{Correct: A} --- it sets the metric name.\\
		\vspace{0.5em}
		\textbf{Wrong:}
		\begin{itemize}
			\item B is \texttt{setLabelCol}.
			\item C is \texttt{setPredictionCol}.
			\item D is unrelated.
		\end{itemize}
	\end{frame}
	
	% ---------------------------------------------------------------
	\begin{frame}{Question 238}
		What does \texttt{setLabelCol()} do on an evaluator?
		\begin{enumerate}[label=\Alph*.]
			\item Sets the target column
			\item Sets the metric
			\item Sets the prediction column
			\item Sets the model
		\end{enumerate}
	\end{frame}
	\begin{frame}{Answer 238}
		\textbf{Correct: A} --- it sets the label column.\\
		\vspace{0.5em}
		\textbf{Wrong:}
		\begin{itemize}
			\item B is \texttt{setMetricName}.
			\item C is \texttt{setPredictionCol}.
			\item D is unrelated.
		\end{itemize}
	\end{frame}
	
	% ---------------------------------------------------------------
	\begin{frame}{Question 239}
		What does \texttt{setPredictionCol()} do on an evaluator?
		\begin{enumerate}[label=\Alph*.]
			\item Sets the column containing predictions
			\item Sets the target column
			\item Sets the metric
			\item Sets the model
		\end{enumerate}
	\end{frame}
	\begin{frame}{Answer 239}
		\textbf{Correct: A} --- it sets the prediction column.\\
		\vspace{0.5em}
		\textbf{Wrong:}
		\begin{itemize}
			\item B is \texttt{setLabelCol}.
			\item C is \texttt{setMetricName}.
			\item D is unrelated.
		\end{itemize}
	\end{frame}
	
	% ---------------------------------------------------------------
	\begin{frame}{Question 240}
		What does \texttt{CrossValidator.fit()} return?
		\begin{enumerate}[label=\Alph*.]
			\item A \texttt{CrossValidatorModel} with the best model
			\item A DataFrame
			\item A metric
			\item A cluster
		\end{enumerate}
	\end{frame}
	\begin{frame}{Answer 240}
		\textbf{Correct: A} --- it returns a \texttt{CrossValidatorModel}.\\
		\vspace{0.5em}
		\textbf{Wrong:}
		\begin{itemize}
			\item B, C, D are not returned.
		\end{itemize}
	\end{frame}
	
	% ---------------------------------------------------------------
	\begin{frame}{Question 241}
		How do you access the best model from a \texttt{CrossValidatorModel}?
		\begin{enumerate}[label=\Alph*.]
			\item \texttt{.bestModel}
			\item \texttt{.model}
			\item \texttt{.best}
			\item \texttt{.getModel()}
		\end{enumerate}
	\end{frame}
	\begin{frame}{Answer 241}
		\textbf{Correct: A} --- \texttt{bestModel} holds the best model.\\
		\vspace{0.5em}
		\textbf{Wrong:}
		\begin{itemize}
			\item B, C, D are not the attribute.
		\end{itemize}
	\end{frame}
	
	% ---------------------------------------------------------------
	\begin{frame}{Question 242}
		What does \texttt{avgMetrics} on a \texttt{CrossValidatorModel} contain?
		\begin{enumerate}[label=\Alph*.]
			\item The average metric for each hyperparameter combination
			\item The metric for one fold
			\item The best model
			\item The feature importance
		\end{enumerate}
	\end{frame}
	\begin{frame}{Answer 242}
		\textbf{Correct: A} --- \texttt{avgMetrics} has per-combination scores.\\
		\vspace{0.5em}
		\textbf{Wrong:}
		\begin{itemize}
			\item B is one fold.
			\item C is \texttt{bestModel}.
			\item D is importance.
		\end{itemize}
	\end{frame}
	
	% ---------------------------------------------------------------
	\begin{frame}{Question 243}
		Why is \texttt{cache()} useful before cross-validation?
		\begin{enumerate}[label=\Alph*.]
			\item It avoids recomputing data across folds
			\item It tunes hyperparameters
			\item It registers models
			\item It deploys models
		\end{enumerate}
	\end{frame}
	\begin{frame}{Answer 243}
		\textbf{Correct: A} --- caching avoids recomputation across folds.\\
		\vspace{0.5em}
		\textbf{Wrong:}
		\begin{itemize}
			\item B, C, D are false.
		\end{itemize}
	\end{frame}
	
	% ---------------------------------------------------------------
	\begin{frame}{Question 244}
		What does \texttt{ParamGridBuilder.addGrid()} do?
		\begin{enumerate}[label=\Alph*.]
			\item Adds a hyperparameter and its values to the grid
			\item Adds a model
			\item Adds a feature
			\item Adds a metric
		\end{enumerate}
	\end{frame}
	\begin{frame}{Answer 244}
		\textbf{Correct: A} --- it adds a parameter grid.\\
		\vspace{0.5em}
		\textbf{Wrong:}
		\begin{itemize}
			\item B, C, D are false.
		\end{itemize}
	\end{frame}
	
	% ---------------------------------------------------------------
	\begin{frame}{Question 245}
		What does \texttt{TrainValidationSplit.setTrainRatio(0.8)} mean?
		\begin{enumerate}[label=\Alph*.]
			\item 80\% of data is used for training
			\item 80\% is used for validation
			\item 80\% is test
			\item 80\% is dropped
		\end{enumerate}
	\end{frame}
	\begin{frame}{Answer 245}
		\textbf{Correct: A} --- 80\% trains, 20\% validates.\\
		\vspace{0.5em}
		\textbf{Wrong:}
		\begin{itemize}
			\item B, C, D misread the ratio.
		\end{itemize}
	\end{frame}
	
	% ---------------------------------------------------------------
	\begin{frame}{Question 246}
		What does \texttt{cvModel.bestModel} refer to when using \texttt{TrainValidationSplit}?
		\begin{enumerate}[label=\Alph*.]
			\item The best model from the split
			\item The metric
			\item The validation data
			\item The training data
		\end{enumerate}
	\end{frame}
	\begin{frame}{Answer 246}
		\textbf{Correct: A} --- \texttt{bestModel} is the best model.\\
		\vspace{0.5em}
		\textbf{Wrong:}
		\begin{itemize}
			\item B, C, D are not models.
		\end{itemize}
	\end{frame}
	
	% ---------------------------------------------------------------
	\begin{frame}{Question 247}
		Which Spark ML classifier predicts probabilities?
		\begin{enumerate}[label=\Alph*.]
			\item \texttt{LogisticRegression}
			\item \texttt{LinearRegression}
			\item \texttt{KMeans}
			\item \texttt{ALS}
		\end{enumerate}
	\end{frame}
	\begin{frame}{Answer 247}
		\textbf{Correct: A} --- logistic regression predicts probabilities.\\
		\vspace{0.5em}
		\textbf{Wrong:}
		\begin{itemize}
			\item B is regression.
			\item C is clustering.
			\item D is recommendation.
		\end{itemize}
	\end{frame}
	
	% ---------------------------------------------------------------
	\begin{frame}{Question 248}
		What does the sigmoid function do in logistic regression?
		\begin{enumerate}[label=\Alph*.]
			\item Maps a linear score to a probability
			\item Scales features
			\item Encodes categories
			\item Imputes values
		\end{enumerate}
	\end{frame}
	\begin{frame}{Answer 248}
		\textbf{Correct: A} --- sigmoid maps scores to probabilities.\\
		\vspace{0.5em}
		\textbf{Wrong:}
		\begin{itemize}
			\item B, C, D are other steps.
		\end{itemize}
	\end{frame}
	
	% ---------------------------------------------------------------
	\begin{frame}{Question 249}
		What is the decision threshold in binary classification?
		\begin{enumerate}[label=\Alph*.]
			\item The probability cutoff for predicting the positive class
			\item The number of trees
			\item The learning rate
			\item The feature count
		\end{enumerate}
	\end{frame}
	\begin{frame}{Answer 249}
		\textbf{Correct: A} --- the threshold converts probabilities to labels.\\
		\vspace{0.5em}
		\textbf{Wrong:}
		\begin{itemize}
			\item B, C, D are unrelated.
		\end{itemize}
	\end{frame}
	
	% ---------------------------------------------------------------
	\begin{frame}{Question 250}
		Why might you lower the decision threshold?
		\begin{enumerate}[label=\Alph*.]
			\item To increase recall
			\item To increase precision
			\item To reduce variance
			\item To reduce training time
		\end{enumerate}
	\end{frame}
	\begin{frame}{Answer 250}
		\textbf{Correct: A} --- lowering the threshold usually raises recall.\\
		\vspace{0.5em}
		\textbf{Wrong:}
		\begin{itemize}
			\item B is the opposite tradeoff.
			\item C, D are unrelated.
		\end{itemize}
	\end{frame}
	
	% ---------------------------------------------------------------
	\begin{frame}{Question 251}
		What does a model's feature importance show?
		\begin{enumerate}[label=\Alph*.]
			\item Which features contribute most to predictions		\item The training time
			\item The cluster size
			\item The data size
		\end{enumerate}
	\end{frame}
	\begin{frame}{Answer 251}
		\textbf{Correct: A} --- importance shows feature contribution.\\
		\vspace{0.5em}
		\textbf{Wrong:}
		\begin{itemize}
			\item B, C, D are unrelated.
		\end{itemize}
	\end{frame}
	
	% ---------------------------------------------------------------
	\begin{frame}{Question 252}
		Which model type supports feature importance in Spark ML?
		\begin{enumerate}[label=\Alph*.]
			\item Tree-based models
			\item KMeans
			\item ALS
			\item PCA
		\end{enumerate}
	\end{frame}
	\begin{frame}{Answer 252}
		\textbf{Correct: A} --- tree models expose feature importance.\\
		\vspace{0.5em}
		\textbf{Wrong:}
		\begin{itemize}
			\item B, C, D do not expose importance.
		\end{itemize}
	\end{frame}
	
	% ---------------------------------------------------------------
	\begin{frame}{Question 253}
		What is a partial dependence plot?
		\begin{enumerate}[label=\Alph*.]
			\item A plot showing the effect of one feature on predictions
			\item A plot of feature importance
			\item A plot of training error
			\item A plot of validation error
		\end{enumerate}
	\end{frame}
	\begin{frame}{Answer 253}
		\textbf{Correct: A} --- partial dependence shows feature effects.\\
		\vspace{0.5em}
		\textbf{Wrong:}
		\begin{itemize}
			\item B, C, D are other plots.
		\end{itemize}
	\end{frame}
	
	% ---------------------------------------------------------------
	\begin{frame}{Question 254}
		What is SHAP used for?
		\begin{enumerate}[label=\Alph*.]
			\item Explaining model predictions with feature attributions
			\item Training models
			\item Deploying models
			\item Storing models
		\end{enumerate}
	\end{frame}
	\begin{frame}{Answer 254}
		\textbf{Correct: A} --- SHAP explains predictions.\\
		\vspace{0.5em}
		\textbf{Wrong:}
		\begin{itemize}
			\item B, C, D are not its purpose.
		\end{itemize}
	\end{frame}
	
	% ---------------------------------------------------------------
	\begin{frame}{Question 255}
		What is LIME used for?
		\begin{enumerate}[label=\Alph*.]
			\item Local explanations of individual predictions
			\item Global training
			\item Model deployment
			\item Data storage
		\end{enumerate}
	\end{frame}
	\begin{frame}{Answer 255}
		\textbf{Correct: A} --- LIME gives local explanations.\\
		\vspace{0.5em}
		\textbf{Wrong:}
		\begin{itemize}
			\item B, C, D are not its purpose.
		\end{itemize}
	\end{frame}
	
	% ---------------------------------------------------------------
	\begin{frame}{Question 256}
		What is a confusion matrix used for?
		\begin{enumerate}[label=\Alph*.]
			\item Evaluating classification outcomes
			\item Evaluating regression outcomes
			\item Clustering
			\item Dimensionality reduction
		\end{enumerate}
	\end{frame}
	\begin{frame}{Answer 256}
		\textbf{Correct: A} --- confusion matrices evaluate classifiers.\\
		\vspace{0.5em}
		\textbf{Wrong:}
		\begin{itemize}
			\item B, C, D use other methods.
		\end{itemize}
	\end{frame}
	
	% ---------------------------------------------------------------
	\begin{frame}{Question 257}
		What does the \texttt{numTrees} parameter control in a random forest?
		\begin{enumerate}[label=\Alph*.]
			\item The number of trees in the forest
			\item The tree depth
			\item The learning rate
			\item The number of features
		\end{enumerate}
	\end{frame}
	\begin{frame}{Answer 257}
		\textbf{Correct: A} --- \texttt{numTrees} sets tree count.\\
		\vspace{0.5em}
		\textbf{Wrong:}
		\begin{itemize}
			\item B is \texttt{maxDepth}.
			\item C is \texttt{stepSize}.
			\item D is \texttt{featureSubsetStrategy}.
		\end{itemize}
	\end{frame}
	
	% ---------------------------------------------------------------
	\begin{frame}{Question 258}
		What does \texttt{maxBins} control in Spark ML trees?
		\begin{enumerate}[label=\Alph*.]
			\item The number of bins for continuous feature splits
			\item The number of trees
			\item The learning rate
			\item The number of features
		\end{enumerate}
	\end{frame}
	\begin{frame}{Answer 258}
		\textbf{Correct: A} --- \texttt{maxBins} controls binning.\\
		\vspace{0.5em}
		\textbf{Wrong:}
		\begin{itemize}
			\item B is \texttt{numTrees}.
			\item C is \texttt{stepSize}.
			\item D is feature strategy.
		\end{itemize}
	\end{frame}
	
	% ---------------------------------------------------------------
	\begin{frame}{Question 259}
		What does \texttt{stepSize} control in gradient boosted trees?
		\begin{enumerate}[label=\Alph*.]
			\item The learning rate
			\item The number of trees
			\item The tree depth
			\item The number of features
		\end{enumerate}
	\end{frame}
	\begin{frame}{Answer 259}
		\textbf{Correct: A} --- \texttt{stepSize} is the learning rate.\\
		\vspace{0.5em}
		\textbf{Wrong:}
		\begin{itemize}
			\item B is \texttt{maxIter}.
			\item C is \texttt{maxDepth}.
			\item D is feature strategy.
		\end{itemize}
	\end{frame}
	
	% ---------------------------------------------------------------
	\begin{frame}{Question 260}
		What does \texttt{subsamplingRate} control in gradient boosted trees?
		\begin{enumerate}[label=\Alph*.]
			\item The fraction of data used for each tree
			\item The number of trees
			\item The tree depth
			\item The number of features
		\end{enumerate}
	\end{frame}
	\begin{frame}{Answer 260}
		\textbf{Correct: A} --- it controls per-tree data sampling.\\
		\vspace{0.5em}
		\textbf{Wrong:}
		\begin{itemize}
			\item B, C, D are other parameters.
		\end{itemize}
	\end{frame}
	
	% ---------------------------------------------------------------
	\begin{frame}{Question 261}
		What does \texttt{featureSubsetStrategy} control in random forests?
		\begin{enumerate}[label=\Alph*.]
			\item How many features are considered per split
			\item The number of trees
			\item The tree depth
			\item The learning rate
		\end{enumerate}
	\end{frame}
	\begin{frame}{Answer 261}
		\textbf{Correct: A} --- it controls features per split.\\
		\vspace{0.5em}
		\textbf{Wrong:}
		\begin{itemize}
			\item B, C, D are other parameters.
		\end{itemize}
	\end{frame}
	
	% ---------------------------------------------------------------
	\begin{frame}{Question 262}
		What does \texttt{minInstancesPerNode} control?
		\begin{enumerate}[label=\Alph*.]
			\item The minimum number of rows required to split a node
			\item The number of trees
			\item The learning rate
			\item The number of features
		\end{enumerate}
	\end{frame}
	\begin{frame}{Answer 262}
		\textbf{Correct: A} --- it sets the minimum rows per node.\\
		\vspace{0.5em}
		\textbf{Wrong:}
		\begin{itemize}
			\item B, C, D are other parameters.
		\end{itemize}
	\end{frame}
	
	% ---------------------------------------------------------------
	\begin{frame}{Question 263}
		What is regularization in linear models?
		\begin{enumerate}[label=\Alph*.]
			\item A penalty on coefficient size to reduce overfitting
			\item A way to add features
			\item A way to remove rows
			\item A clustering method
		\end{enumerate}
	\end{frame}
	\begin{frame}{Answer 263}
		\textbf{Correct: A} --- regularization penalizes coefficients.\\
		\vspace{0.5em}
		\textbf{Wrong:}
		\begin{itemize}
			\item B, C, D are false.
		\end{itemize}
	\end{frame}
	
	% ---------------------------------------------------------------
	\begin{frame}{Question 264}
		What is the difference between L1 and L2 regularization?
		\begin{enumerate}[label=\Alph*.]
			\item L1 can zero out coefficients; L2 shrinks them
			\item L2 zeros coefficients; L1 shrinks them
			\item They are identical
			\item Neither affects coefficients
		\end{enumerate}
	\end{frame}
	\begin{frame}{Answer 264}
		\textbf{Correct: A} --- L1 is sparse; L2 shrinks.\\
		\vspace{0.5em}
		\textbf{Wrong:}
		\begin{itemize}
			\item B is reversed.
			\item C, D are false.
		\end{itemize}
	\end{frame}
	
	% ---------------------------------------------------------------
	\begin{frame}{Question 265}
		What does \texttt{regParam} control in Spark ML linear models?
		\begin{enumerate}[label=\Alph*.]
			\item The regularization strength
			\item The learning rate
			\item The number of iterations
			\item The number of features
		\end{enumerate}
	\end{frame}
	\begin{frame}{Answer 265}
		\textbf{Correct: A} --- \texttt{regParam} sets regularization.\\
		\vspace{0.5em}
		\textbf{Wrong:}
		\begin{itemize}
			\item B is \texttt{stepSize}.
			\item C is \texttt{maxIter}.
			\item D is feature count.
		\end{itemize}
	\end{frame}
	
	% ---------------------------------------------------------------
	\begin{frame}{Question 266}
		What does \texttt{elasticNetParam} control?
		\begin{enumerate}[label=\Alph*.]
			\item The mix of L1 and L2 regularization
			\item The number of trees
			\item The number of iterations
			\item The feature count
		\end{enumerate}
	\end{frame}
	\begin{frame}{Answer 266}
		\textbf{Correct: A} --- it controls the L1/L2 mix.\\
		\vspace{0.5em}
		\textbf{Wrong:}
		\begin{itemize}
			\item B, C, D are other parameters.
		\end{itemize}
	\end{frame}
	
	% ---------------------------------------------------------------
	\begin{frame}{Question 267}
		What is a sparse vector in Spark ML?
		\begin{enumerate}[label=\Alph*.]
			\item A vector with mostly zero values
			\item A vector with all nonzero values
			\item A vector of strings
			\item A vector of labels
		\end{enumerate}
	\end{frame}
	\begin{frame}{Answer 267}
		\textbf{Correct: A} --- sparse vectors are mostly zeros.\\
		\vspace{0.5em}
		\textbf{Wrong:}
		\begin{itemize}
			\item B is dense.
			\item C, D are other data types.
		\end{itemize}
	\end{frame}
	
	% ---------------------------------------------------------------
	\begin{frame}{Question 268}
		Why does one-hot encoding produce sparse vectors?
		\begin{enumerate}[label=\Alph*.]
			\item Most categories are zero for any given row
			\item All values are nonzero
			\item It encodes numeric data
			\item It scales features
		\end{enumerate}
	\end{frame}
	\begin{frame}{Answer 268}
		\textbf{Correct: A} --- one-hot vectors are mostly zeros.\\
		\vspace{0.5em}
		\textbf{Wrong:}
		\begin{itemize}
			\item B, C, D are false.
		\end{itemize}
	\end{frame}
	
	% ---------------------------------------------------------------
	\begin{frame}{Question 269}
		What does \texttt{VectorSlicer} do?
		\begin{enumerate}[label=\Alph*.]
			\item Selects a subset of features from a vector
			\item Scales a vector
			\item Encodes a vector
			\item Imputes a vector
		\end{enumerate}
	\end{frame}
	\begin{frame}{Answer 269}
		\textbf{Correct: A} --- VectorSlicer selects features.\\
		\vspace{0.5em}
		\textbf{Wrong:}
		\begin{itemize}
			\item B, C, D are other transformers.
		\end{itemize}
	\end{frame}
	
	% ---------------------------------------------------------------
	\begin{frame}{Question 270}
		What is the purpose of a \texttt{SQLTransformer}?
		\begin{enumerate}[label=\Alph*.]
			\item It applies a SQL query as a transformation
			\item It trains a model
			\item It registers a model
			\item It deploys a model
		\end{enumerate}
	\end{frame}
	\begin{frame}{Answer 270}
		\textbf{Correct: A} --- SQLTransformer applies SQL.\\
		\vspace{0.5em}
		\textbf{Wrong:}
		\begin{itemize}
			\item B, C, D are false.
		\end{itemize}
	\end{frame}
	
	% ---------------------------------------------------------------
	\begin{frame}{Question 271}
		What is a Delta table?
		\begin{enumerate}[label=\Alph*.]
			\item A table stored in Delta Lake with ACID transactions
			\item A temporary view
			\item A Python list
			\item A model
		\end{enumerate}
	\end{frame}
	\begin{frame}{Answer 271}
		\textbf{Correct: A} --- Delta tables provide ACID storage.\\
		\vspace{0.5em}
		\textbf{Wrong:}
		\begin{itemize}
			\item B, C, D are not Delta tables.
		\end{itemize}
	\end{frame}
	
	% ---------------------------------------------------------------
	\begin{frame}{Question 272}
		What is time travel in Delta Lake?
		\begin{enumerate}[label=\Alph*.]
			\item Querying previous versions of a table
			\item Time-based clustering
			\item A scheduling feature
			\item A model type
		\end{enumerate}
	\end{frame}
	\begin{frame}{Answer 272}
		\textbf{Correct: A} --- time travel queries old versions.\\
		\vspace{0.5em}
		\textbf{Wrong:}
		\begin{itemize}
			\item B, C, D are false.
		\end{itemize}
	\end{frame}
	
	% ---------------------------------------------------------------
	\begin{frame}{Question 273}
		What does \texttt{OPTIMIZE} do on a Delta table?
		\begin{enumerate}[label=\Alph*.]
			\item Compacts small files
			\item Deletes the table
			\item Adds a column
			\item Changes the schema
		\end{enumerate}
	\end{frame}
	\begin{frame}{Answer 273}
		\textbf{Correct: A} --- OPTIMIZE compacts files.\\
		\vspace{0.5em}
		\textbf{Wrong:}
		\begin{itemize}
			\item B, C, D are false.
		\end{itemize}
	\end{frame}
	
	% ---------------------------------------------------------------
	\begin{frame}{Question 274}
		What is \texttt{VACUUM} used for on a Delta table?
		\begin{enumerate}[label=\Alph*.]
			\item Removing old files not needed for time travel
			\item Adding files
			\item Compacting files
			\item Changing the schema
		\end{enumerate}
	\end{frame}
	\begin{frame}{Answer 274}
		\textbf{Correct: A} --- VACUUM removes stale files.\\
		\vspace{0.5em}
		\textbf{Wrong:}
		\begin{itemize}
			\item B, C, D are false.
		\end{itemize}
	\end{frame}
	
	% ---------------------------------------------------------------
	\begin{frame}{Question 275}
		Which Unity Catalog object contains schemas?
		\begin{enumerate}[label=\Alph*.]
			\item Catalog
			\item Table
			\item View
			\item Volume
		\end{enumerate}
	\end{frame}
	\begin{frame}{Answer 275}
		\textbf{Correct: A} --- catalogs contain schemas.\\
		\vspace{0.5em}
		\textbf{Wrong:}
		\begin{itemize}
			\item B, C, D are lower-level objects.
		\end{itemize}
	\end{frame}
	
	% ---------------------------------------------------------------
	\begin{frame}{Question 276}
		Which Unity Catalog object contains tables?
		\begin{enumerate}[label=\Alph*.]
			\item Schema
			\item Catalog
			\item Volume
			\item Workspace
		\end{enumerate}
	\end{frame}
	\begin{frame}{Answer 276}
		\textbf{Correct: A} --- schemas contain tables.\\
		\vspace{0.5em}
		\textbf{Wrong:}
		\begin{itemize}
			\item B contains schemas.
			\item C is for files.
			\item D is not a UC object.
		\end{itemize}
	\end{frame}
	
	% ---------------------------------------------------------------
	\begin{frame}{Question 277}
		What is a Unity Catalog volume used for?
		\begin{enumerate}[label=\Alph*.]
			\item Storing files under governance
			\item Storing tables
			\item Storing models
			\item Storing functions
		\end{enumerate}
	\end{frame}
	\begin{frame}{Answer 277}
		\textbf{Correct: A} --- volumes store files.\\
		\vspace{0.5em}
		\textbf{Wrong:}
		\begin{itemize}
			\item B is tables.
			\item C is models.
			\item D is functions.
		\end{itemize}
	\end{frame}
	
	% ---------------------------------------------------------------
	\begin{frame}{Question 278}
		Which privilege allows reading a table in Unity Catalog?
		\begin{enumerate}[label=\Alph*.]
			\item \texttt{SELECT}
			\item \texttt{READ}
			\item \texttt{VIEW}
			\item \texttt{GET}
		\end{enumerate}
	\end{frame}
	\begin{frame}{Answer 278}
		\textbf{Correct: A} --- \texttt{SELECT} grants read access.\\
		\vspace{0.5em}
		\textbf{Wrong:}
		\begin{itemize}
			\item B, C, D are not the privilege names.
		\end{itemize}
	\end{frame}
	
	% ---------------------------------------------------------------
	\begin{frame}{Question 279}
		Which privilege allows writing to a table in Unity Catalog?
		\begin{enumerate}[label=\Alph*.]
			\item \texttt{MODIFY}
			\item \texttt{WRITE}
			\item \texttt{INSERT}
			\item \texttt{UPDATE}
		\end{enumerate}
	\end{frame}
	\begin{frame}{Answer 279}
		\textbf{Correct: A} --- \texttt{MODIFY} grants write access.\\
		\vspace{0.5em}
		\textbf{Wrong:}
		\begin{itemize}
			\item B, C, D are not the privilege names.
		\end{itemize}
	\end{frame}
	
	% ---------------------------------------------------------------
	\begin{frame}{Question 280}
		What is a Databricks Job?
		\begin{enumerate}[label=\Alph*.]
			\item A scheduled or triggered workflow of tasks
			\item A cluster
			\item A notebook
			\item A model
		\end{enumerate}
	\end{frame}
	\begin{frame}{Answer 280}
		\textbf{Correct: A} --- Jobs run tasks on a schedule or trigger.\\
		\vspace{0.5em}
		\textbf{Wrong:}
		\begin{itemize}
			\item B, C, D are not jobs.
		\end{itemize}
	\end{frame}
	
	% ---------------------------------------------------------------
	\begin{frame}{Question 281}
		What is the difference between all-purpose and job clusters?
		\begin{enumerate}[label=\Alph*.]
			\item All-purpose clusters are interactive; job clusters are for automated jobs
			\item They are identical
			\item Job clusters are interactive
			\item All-purpose clusters are only for SQL
		\end{enumerate}
	\end{frame}
	\begin{frame}{Answer 281}
		\textbf{Correct: A} --- interactive vs.\ job compute.\\
		\vspace{0.5em}
		\textbf{Wrong:}
		\begin{itemize}
			\item B, C, D are false.
		\end{itemize}
	\end{frame}
	
	% ---------------------------------------------------------------
	\begin{frame}{Question 282}
		What does a Databricks Repo integrate with?
		\begin{enumerate}[label=\Alph*.]
			\item A Git provider
			\item A SQL warehouse
			\item A model registry
			\item A feature store
		\end{enumerate}
	\end{frame}
	\begin{frame}{Answer 282}
		\textbf{Correct: A} --- Repos integrate with Git.\\
		\vspace{0.5em}
		\textbf{Wrong:}
		\begin{itemize}
			\item B, C, D are other services.
		\end{itemize}
	\end{frame}
	
	% ---------------------------------------------------------------
	\begin{frame}{Question 283}
		Why use a Git branch for feature development in Databricks Repos?
		\begin{enumerate}[label=\Alph*.]
			\item To isolate changes from production code
			\item To train faster
			\item To save storage
			\item To improve accuracy
		\end{enumerate}
	\end{frame}
	\begin{frame}{Answer 283}
		\textbf{Correct: A} --- branches isolate development.\\
		\vspace{0.5em}
		\textbf{Wrong:}
		\begin{itemize}
			\item B, C, D are false.
		\end{itemize}
	\end{frame}
	
	% ---------------------------------------------------------------
	\begin{frame}{Question 284}
		What is the matrix view in a Databricks Job run used for?
		\begin{enumerate}[label=\Alph*.]
			\item Visualizing task statuses and failures
			\item Training models
			\item Registering models
			\item Writing SQL
		\end{enumerate}
	\end{frame}
	\begin{frame}{Answer 284}
		\textbf{Correct: A} --- the matrix view shows task status.\\
		\vspace{0.5em}
		\textbf{Wrong:}
		\begin{itemize}
			\item B, C, D are unrelated.
		\end{itemize}
	\end{frame}
	
	% ---------------------------------------------------------------
	\begin{frame}{Question 285}
		What is a notebook widget in Databricks?
		\begin{enumerate}[label=\Alph*.]
			\item A parameter input for a notebook
			\item A model
			\item A cluster
			\item A table
		\end{enumerate}
	\end{frame}
	\begin{frame}{Answer 285}
		\textbf{Correct: A} --- widgets are notebook parameters.\\
		\vspace{0.5em}
		\textbf{Wrong:}
		\begin{itemize}
			\item B, C, D are not widgets.
		\end{itemize}
	\end{frame}
	
	% ---------------------------------------------------------------
	\begin{frame}{Question 286}
		What does \texttt{dbutils.widgets.text()} do?
		\begin{enumerate}[label=\Alph*.]
			\item Creates a text input widget
			\item Creates a cluster
			\item Creates a table
			\item Creates a model
		\end{enumerate}
	\end{frame}
	\begin{frame}{Answer 286}
		\textbf{Correct: A} --- it creates a text widget.\\
		\vspace{0.5em}
		\textbf{Wrong:}
		\begin{itemize}
			\item B, C, D are false.
		\end{itemize}
	\end{frame}
	
	% ---------------------------------------------------------------
	\begin{frame}{Question 287}
		What does \texttt{dbutils.secrets.get()} do?
		\begin{enumerate}[label=\Alph*.]
			\item Retrieves a secret value securely
			\item Creates a secret scope
			\item Deletes a secret
			\item Lists secrets
		\end{enumerate}
	\end{frame}
	\begin{frame}{Answer 287}
		\textbf{Correct: A} --- it retrieves a secret.\\
		\vspace{0.5em}
		\textbf{Wrong:}
		\begin{itemize}
			\item B, C, D are other secret operations.
		\end{itemize}
	\end{frame}
	
	% ---------------------------------------------------------------
	\begin{frame}{Question 288}
		Why should secrets not be hard-coded in notebooks?
		\begin{enumerate}[label=\Alph*.]
			\item They can be exposed in version control and logs
			\item They slow training
			\item They reduce accuracy
			\item They increase storage
		\end{enumerate}
	\end{frame}
	\begin{frame}{Answer 288}
		\textbf{Correct: A} --- hard-coded secrets leak.\\
		\vspace{0.5em}
		\textbf{Wrong:}
		\begin{itemize}
			\item B, C, D are false.
		\end{itemize}
	\end{frame}
	
	% ---------------------------------------------------------------
	\begin{frame}{Question 289}
		Which serving endpoint type supports low-latency REST inference?
		\begin{enumerate}[label=\Alph*.]
			\item Real-time serving
			\item Batch serving
			\item Streaming serving
			\item Notebook serving
		\end{enumerate}
	\end{frame}
	\begin{frame}{Answer 289}
		\textbf{Correct: A} --- real-time serving is low-latency REST.\\
		\vspace{0.5em}
		\textbf{Wrong:}
		\begin{itemize}
			\item B, C are other modes.
			\item D is not a serving mode.
		\end{itemize}
	\end{frame}
	
	% ---------------------------------------------------------------
	\begin{frame}{Question 290}
		What does a serving endpoint require to route to multiple model versions?
		\begin{enumerate}[label=\Alph*.]
			\item Traffic splitting configuration
			\item A new cluster
			\item A new table
			\item A new notebook
		\end{enumerate}
	\end{frame}
	\begin{frame}{Answer 290}
		\textbf{Correct: A} --- traffic splitting enables multi-version routing.\\
		\vspace{0.5em}
		\textbf{Wrong:}
		\begin{itemize}
			\item B, C, D are unrelated.
		\end{itemize}
	\end{frame}
	
	% ---------------------------------------------------------------
	\begin{frame}{Question 291}
		What is the purpose of authentication on a serving endpoint?
		\begin{enumerate}[label=\Alph*.]
			\item To control who can call the endpoint
			\item To train the model
			\item To store the model
			\item To register the model
		\end{enumerate}
	\end{frame}
	\begin{frame}{Answer 291}
		\textbf{Correct: A} --- auth controls access.\\
		\vspace{0.5em}
		\textbf{Wrong:}
		\begin{itemize}
			\item B, C, D are unrelated.
		\end{itemize}
	\end{frame}
	
	% ---------------------------------------------------------------
	\begin{frame}{Question 292}
		What is a batch inference job typically scheduled with?
		\begin{enumerate}[label=\Alph*.]
			\item A Databricks Job
			\item A SQL warehouse
			\item A notebook widget
			\item A Repo
		\end{enumerate}
	\end{frame}
	\begin{frame}{Answer 292}
		\textbf{Correct: A} --- Jobs schedule batch inference.\\
		\vspace{0.5em}
		\textbf{Wrong:}
		\begin{itemize}
			\item B is SQL compute.
			\item C is UI.
			\item D is Git.
		\end{itemize}
	\end{frame}
	
	% ---------------------------------------------------------------
	\begin{frame}{Question 293}
		What is the main risk of training on stale features?
		\begin{enumerate}[label=\Alph*.]
			\item The model learns patterns that do not hold in production
			\item Training becomes faster
			\item Accuracy always improves
			\item Storage increases
		\end{enumerate}
	\end{frame}
	\begin{frame}{Answer 293}
		\textbf{Correct: A} --- stale features cause misalignment.\\
		\vspace{0.5em}
		\textbf{Wrong:}
		\begin{itemize}
			\item B, C, D are false.
		\end{itemize}
	\end{frame}
	
	% ---------------------------------------------------------------
	\begin{frame}{Question 294}
		Why is feature lineage important?
		\begin{enumerate}[label=\Alph*.]
			\item It shows where features come from and how they were computed
			\item It speeds training
			\item It reduces storage
			\item It improves accuracy
		\end{enumerate}
	\end{frame}
	\begin{frame}{Answer 294}
		\textbf{Correct: A} --- lineage tracks feature origin.\\
		\vspace{0.5em}
		\textbf{Wrong:}
		\begin{itemize}
			\item B, C, D are false.
		\end{itemize}
	\end{frame}
	
	% ---------------------------------------------------------------
	\begin{frame}{Question 295}
		What is a model signature used for in serving?
		\begin{enumerate}[label=\Alph*.]
			\item Validating input and output schemas
			\item Training the model
			\item Registering the model
			\item Deleting the model
		\end{enumerate}
	\end{frame}
	\begin{frame}{Answer 295}
		\textbf{Correct: A} --- signatures validate schemas.\\
		\vspace{0.5em}
		\textbf{Wrong:}
		\begin{itemize}
			\item B, C, D are unrelated.
		\end{itemize}
	\end{frame}
	
	% ---------------------------------------------------------------
	\begin{frame}{Question 296}
		What is the purpose of a data quality check before training?
		\begin{enumerate}[label=\Alph*.]
			\item To catch missing, invalid, or drifted values
			\item To speed training
			\item To register models
			\item To deploy models
		\end{enumerate}
	\end{frame}
	\begin{frame}{Answer 296}
		\textbf{Correct: A} --- quality checks catch bad data.\\
		\vspace{0.5em}
		\textbf{Wrong:}
		\begin{itemize}
			\item B, C, D are unrelated.
		\end{itemize}
	\end{frame}
	
	% ---------------------------------------------------------------
	\begin{frame}{Question 297}
		Why is reproducibility important in ML?
		\begin{enumerate}[label=\Alph*.]
			\item It lets others verify and reproduce results
			\item It speeds training
			\item It reduces storage
			\item It improves accuracy
		\end{enumerate}
	\end{frame}
	\begin{frame}{Answer 297}
		\textbf{Correct: A} --- reproducibility supports verification.\\
		\vspace{0.5em}
		\textbf{Wrong:}
		\begin{itemize}
			\item B, C, D are not the main reason.
		\end{itemize}
	\end{frame}
	
	% ---------------------------------------------------------------
	\begin{frame}{Question 298}
		What is a model card?
		\begin{enumerate}[label=\Alph*.]
			\item A document describing a model's use, data, performance, and limitations
			\item A model file
			\item A cluster
			\item A notebook
		\end{enumerate}
	\end{frame}
	\begin{frame}{Answer 298}
		\textbf{Correct: A} --- model cards document models.\\
		\vspace{0.5em}
		\textbf{Wrong:}
		\begin{itemize}
			\item B, C, D are not model cards.
		\end{itemize}
	\end{frame}
	
	% ---------------------------------------------------------------
	\begin{frame}{Question 299}
		What is the purpose of monitoring feature drift in production?
		\begin{enumerate}[label=\Alph*.]
			\item To detect when input distributions change
			\item To train the model
			\item To register the model
			\item To delete the model
		\end{enumerate}
	\end{frame}
	\begin{frame}{Answer 299}
		\textbf{Correct: A} --- drift monitoring detects input changes.\\
		\vspace{0.5em}
		\textbf{Wrong:}
		\begin{itemize}
			\item B, C, D are unrelated.
		\end{itemize}
	\end{frame}
	
	% ---------------------------------------------------------------
	\begin{frame}{Question 300}
		What is the best practice for logging secrets in MLflow?
		\begin{enumerate}[label=\Alph*.]
			\item Never log secrets; use secret scopes
			\item Log them as params
			\item Log them as metrics
			\item Log them as artifacts
		\end{enumerate}
	\end{frame}
	\begin{frame}{Answer 300}
		\textbf{Correct: A} --- secrets must not be logged.\\
		\vspace{0.5em}
		\textbf{Wrong:}
		\begin{itemize}
			\item B, C, D expose secrets.
		\end{itemize}
	\end{frame}
	
\end{document}